\begin{trapbox}
	\n
	\begin{description}[style=nextline, leftmargin=2.8em, labelsep=0.5em, itemsep=0.35em]
		\n
		\item[\textbf{\textcolor{CTrap}{易错点一（忽略二次项系）：}}]
			当二次项系数 $a$ 可能为零时，未分类讨论，直接使用判别式判断交点个数。例如在含参二次函数中，若未要求 $a\neq0$，则方程可能退化为一元一次方程，此时交点个数需单独判断。\n
		\item[\textbf{\textcolor{CTrap}{正确理解（确认a≠0）：}}]
			必须首先确认 $a\neq0$ 才能使用判别式结论。若 $a=0$，则原函数不是二次函数，联立后为一次方程，交点个数由一次方程决定。\n
		\item[\textbf{\textcolor{CTrap}{错因分析（条件缺失）：}}]
			做题时未仔细检查题目是否明确 $a\neq0$，或粗心默认 $a\neq0$ 而未考虑参数影响。\n
	\end{description}
	\n\n

	\medskip
	\n
	\begin{description}[style=nextline, leftmargin=2.8em, labelsep=0.5em, itemsep=0.35em]
		\n
		\item[\textbf{\textcolor{CTrap}{易错点二（符号记反）：}}]
			将判别式符号与交点个数的对应关系记反，如认为 $\Delta>0$ 时无交点，$\Delta<0$ 时有两个交点。\n
		\item[\textbf{\textcolor{CTrap}{正确理解（牢固记忆）：}}]
			正确对应：$\Delta>0$ 有两个交点，$\Delta=0$ 有一个交点（相切），$\Delta<0$ 无交点。可借助二次方程根的情况记忆。\n
		\item[\textbf{\textcolor{CTrap}{错因分析（记忆混乱）：}}]
			与二次函数其他判别（如最值判别）混淆，或受错误口诀影响。需反复对比巩固。\n
	\end{description}
	\n\n

	\medskip
	\n
	\begin{description}[style=nextline, leftmargin=2.8em, labelsep=0.5em, itemsep=0.35em]
		\n
		\item[\textbf{\textcolor{CTrap}{易错点三（移项符号错）：}}]
			联立 $y=ax^{2}+bx+c$ 与 $y=kx+m$ 后，移项时出现符号错误，例如写成 $ax^{2}+(b+k)x+(c+m)=0$。\n
		\item[\textbf{\textcolor{CTrap}{正确理解（正确移项）：}}]
			正确移项：$ax^{2}+bx+c - (kx+m)=0$，即 $ax^{2}+(b-k)x+(c-m)=0$。\n
		\item[\textbf{\textcolor{CTrap}{错因分析（运算疏忽）：}}]
			忘记括号，忽略 $kx+m$ 应整体变号，尤其是常数项 $m$ 变为 $-m$。\n
	\end{description}
	\n
\end{trapbox}
